\begin{table}[ht]
\centering
\caption{Selected diplomatic snap episodes in the UN-voting layer. Category entries are posterior mean snap probabilities for the focal actors in that category.}
\label{tab:episode-typology}
\scriptsize
\begin{tabularx}{\textwidth}{@{}r r r r r r r X@{}}
\toprule
\textbf{Year} & \textbf{Index} & \textbf{U.S.} & \textbf{Allies} & \textbf{Rivals} & \textbf{Hedgers} & \textbf{$>0.5$} & \textbf{Signature} \\
\midrule
1991 & 0.406 & 0.110 & 0.431 & 0.408 & 0.400 & 12 & post-Cold War adjustment; allies, rivals, and hedgers move more than the United States \\
1998 & 0.511 & 0.627 & 0.399 & 0.681 & 0.503 & 17 & late-1990s rival and hedger re-coordination, with Turkey, China, Belarus, and Japan prominent \\
2001 & 0.465 & 0.457 & 0.394 & 0.502 & 0.525 & 10 & post-9/11 and Bush-era re-coordination; rivals and hedgers are most exposed \\
2009 & 0.619 & 0.843 & 0.678 & 0.531 & 0.595 & 24 & financial-crisis-era diplomatic re-coordination; broad but less synchronized than 2025 \\
2013 & 0.593 & 0.630 & 0.737 & 0.474 & 0.504 & 22 & pre-Crimea and Arab Spring-era bloc repositioning; allies and rivals are high \\
2019 & 0.498 & 0.505 & 0.526 & 0.512 & 0.451 & 18 & Trump I-era re-coordination; broad but moderate relative to 2025 \\
2023 & 0.544 & 0.740 & 0.585 & 0.446 & 0.554 & 19 & Russia-invasion aftermath; United States, allies, and hedgers are high \\
2025 & 0.956 & 1.000 & 1.000 & 0.896 & 0.944 & 34 & bloc-preserving hegemonic exit; all focal categories snap together \\
\bottomrule
\end{tabularx}
\end{table}
